% Content of the Vlasov blueprint.
% The LeanArchitect-extracted nodes are defined in the aggregate library file below
% (it provides \newleannode/\inputleannode and the per-node \input{...} bodies).
% We then arrange the nodes into a narrative with \inputleannode.

\input{../../.lake/build/blueprint/library/BlueprintDemo}

This blueprint is a reader's map of a Lean~4 / Mathlib formalization of the mean-field
derivation of the Vlasov equation: forward well-posedness, Dobrushin's 1979 stability
estimate, the mean-field limit, and a short-window superposition principle (weak solutions
are Lagrangian). The statements on the cards below are extracted from the elaborated Lean
proof terms, and the edges of the
\href{https://hydrodynamical.github.io/Vlasov_Meanfield_Formalization/dep_graph_document.html}{dependency
graph} are computed from the proofs themselves, not asserted. Where a card corresponds to a
numbered statement of the companion paper, the paper's number appears in its title. The
three apex theorems --- well-posedness, Dobrushin stability, and the superposition
principle --- each carry the axiom footprint
$[\mathtt{propext}, \mathtt{Classical.choice}, \mathtt{Quot.sound}]$: Lean's three standard
axioms, no \texttt{sorry}, no project-specific axioms.

Companion artifacts:
the \href{https://github.com/Hydrodynamical/Vlasov_Meanfield_Formalization/blob/main/PAPER.pdf}{paper}
(\emph{Mathematician in the Loop}),
the \href{https://github.com/Hydrodynamical/Vlasov_Meanfield_Formalization}{repository},
and the \href{https://hydrodynamical.github.io/Vlasov_Meanfield_Formalization/docs/}{API
documentation} --- every declaration, searchable, with source links to the exact commit.
Each card's ``Lean declarations'' button resolves to its API page.

\chapter{From Newton to the Vlasov equation}

The derivation opens with $N$ identical particles in $\mathbb{R}^d$ under the mean-field
force $-\frac{1}{N}\sum_j \nabla W$. The empirical measure of the particle system is not
merely close to a solution of the limit equation: under the standing evenness assumption
it \emph{is} one, exactly, for every $N$. The chapter records the standing assumption, the
particle system, the computation that exposes the $O(1/N)$ diagonal remainder, the
cancellation that kills it, and the two solution concepts --- weak and Lagrangian --- whose
relation the rest of the development settles.

\inputleannode{ass:W}
\inputleannode{def:hamiltonian}
\inputleannode{def:newton}
\inputleannode{def:empirical}
\inputleannode{def:empirical-curve}
\inputleannode{prop:weak-evolution}
\inputleannode{cor:empirical-vlasov}
\inputleannode{def:convolve}
\inputleannode{def:vlasov-sol}
\inputleannode{def:char-flow}
\inputleannode{def:lagrangian-sol}

\chapter[The Kantorovich--Rubinstein W1 distance]{The Kantorovich--Rubinstein $W_1$ distance}

Solutions depend stably on their data, and the right distance is the Wasserstein-$1$
metric. The development defines $W_1$ by its dual sup-formula (Definition~1.4 of the
paper), which makes the non-expansion estimates --- the workhorses of the dynamical
chapters --- immediate. The primal (coupling) face, on which Dobrushin's argument runs, is
recovered by Kantorovich--Rubinstein duality: the hard direction is built from finite
transport problems by cone separation, with no compactness of the ambient space and no
attained optimum anywhere --- every bound in the development is $\varepsilon$-optimal.

\inputleannode{ot:wcost}
\inputleannode{ot:wcost-self}
\inputleannode{ot:wcost-comm}
\inputleannode{ot:wcost-tri}
\inputleannode{ot:wcost-lip}
\inputleannode{ot:w1}
\inputleannode{ot:w1-self}
\inputleannode{ot:w1-comm}
\inputleannode{ot:w1-tri}
\inputleannode{ot:w1-lip}
\inputleannode{ot:w1-fin}
\inputleannode{ot:coupling}
\inputleannode{ot:wcost-coupling}
\inputleannode{ot:w1-coupling}
\inputleannode{ot:w1-le-coupling}
\inputleannode{ot:coupling-le-dual}
\inputleannode{ot:w1-eq-coupling}

\chapter{The characteristic flow and well-posedness}

The Vlasov equation transports mass along the characteristic system
$\dot X = V$, $\dot V = -(\nabla W * \rho_t)(X)$. Existence is a fixed-point argument in
the space of measure curves: freeze the field, solve the frozen linear problem by
transporting the datum along its flow, and contract in the $W_1$-sup metric on a short
window. Windows of a fixed length depending only on $L$ then glue to arbitrary horizons,
so no smallness of the interaction is required --- Theorem~1.3 holds at every Lipschitz
constant.

\inputleannode{def:vlasov-field}
\inputleannode{def:vlasov-sol-on}
\inputleannode{def:lagrangian-sol-on}
\inputleannode{thm:vlasov-wp}

\chapter{Dobrushin stability and the mean-field limit}

Dobrushin's 1979 theorem: two Lagrangian solutions separate at most exponentially in
$W_1$, at the explicit rate $C = 2\max(1,L)$. The proof is his own coupling argument,
followed piece for piece --- couple the data optimally, push the coupling forward along the
pair of flows, bound the force difference by the coupling cost, and close with a Gronwall
inequality in mild form. Applied with one solution the empirical measure of the Newton
dynamics --- an exact weak solution by Chapter~1 --- the estimate yields the mean-field
limit: propagation of chaos in quantitative $W_1$ form.

\inputleannode{thm:dobrushin}
\inputleannode{def:dobrushin-estimate}
\inputleannode{cor:mean-field-limit}

\chapter[The superposition principle: weak implies Lagrangian]{The superposition principle: weak $\Rightarrow$ Lagrangian}

Uniqueness in Theorem~1.3 is uniqueness among \emph{Lagrangian} solutions. Ruling out
phantom weak solutions --- distributional solutions not transported by any flow --- is the
superposition principle, proved here on a short window at the cost of one degree of
regularity (Assumption~2). The argument freezes the field at the solution's own marginal,
establishes the $C^1$ dependence of the flow on its initial point via the variational
equation, enlarges the test class from $C^\infty_c$ to $C^1_c$ to admit transported tests,
and closes with a diagonal chain rule: differentiating $\int \psi_s\,\mathrm{d}f_s$ through
both moving arguments, the transport identity and the weak PDE cancel exactly.

\inputleannode{ass:W2}
\inputleannode{def:fundamental-matrix}
\inputleannode{thm:charflow-fderiv-fundamental}
\inputleannode{thm:variational-eq}
\inputleannode{thm:charflow-inverse-joint}
\inputleannode{thm:test-c1c}
\inputleannode{thm:transport-identity}
\inputleannode{thm:diagonal-chain-rule}
\inputleannode{thm:dual-core}
\inputleannode{thm:weak-lagrangian}
